\paragraph{\trPicard--\trLefschetz{}数据。}\label{sub:picard-lefschetz-data}
我们再次考虑一个 $(n+1)$ 维的\trSmooth{}\trProjectiveVariety，但这一次它显式地\trEmbedding{}于\trProjectiveSpace{}之中：$\bar{X}\subset\CP^N$。如前那样，它应满足
\begin{equation}\label{eq:canonical-bundle-embedding}
\mathcal{K}_{\bar{X}}\cong\scr{O}_{\bar{X}}(-d)
\end{equation}
对某个 $d\in\ZZ$。我们将假设这一\trProjective{}\trEmbedding{}相对于标准\trHomogeneousCoordinates{}$(y_0:\cdots:y_N)$ 处于\textbf{\trGeneralPosition}（general position）。确切的条件将在下面列出，但要点在于：在施加 $\mathrm{PGL}_{N+1}$ 的某个适当元素之后——\trEquivalence{}地，在换用另一套\trHomogeneousCoordinates{}系之后——它们总会成立。

取无穷远\trHyperplane{}为 $\CP^{N-1}=\{y_N=0\}$，把它的补集与 $\CC^N$ 等同起来，坐标为 $(x_0=y_0/y_N,\ldots,x_{N-1}=y_{N-1}/y_N)$。对 $0\le k\le n+1$，取子空间 $H^k=\{x_0=\cdots=x_{k-1}=0\}\subset\CC^N$ 及其闭包 $\bar{H}^k\subset\CP^N$。定义
\begin{equation*}
\begin{aligned}
\bar{X}^k&=\bar{X}\cap\bar{H}^k,\\
X^k&=X\cap H^k,\\
\Psi^k&=x_k\colon X^k\longrightarrow\CC,\\
\varpi^k&=(x_k,x_{k+1})\colon X^k\longrightarrow\CC^2.
\end{aligned}
\end{equation*}
\trGeneric{}地，下面这些将会成立：
\begin{itemize}
\item （对 $0\le k\le n+1$）$\bar{X}$ \trTransverse{}于所有的 $\bar{H}^k$ 与 $\bar{H}^k\cap\CP^{N-1}$。这表明每块 $X^k$ 都是\trSmooth{}的，并且在无穷远处也\trSmooth。
\end{itemize}
注意，按定义 $X^{k+1}=(\Psi^k)^{-1}(0)$，且 $\Psi^{k+1}=\varpi^k|X^{k+1}$。此外，从 \eqref{eq:canonical-bundle-embedding} 我们得到 $\mathcal{K}_{\bar{X}^k}\cong\scr{O}_{\bar{X}^k}(k-d)$。继续罗列这份\trGeneric{}性质清单：
\begin{itemize}
\item （对 $0\le k\le n$）函数 $\Psi^k$ 都只有\trNondegenerate{}的\trCriticalPoint，并且没有两个位于同一根\trFiber{}之中。
\item （对 $0\le k\le n-1$）$D\varpi^k$ 是\trGeneric{}的（即\trTransverse{}于 $\mathrm{Hom}(TX^k,\CC^2)$ 中由低秩\trLinear{}映射组成的那些子簇）。如果 $w\in\CC$ 是 $\Psi^k$ 的一个\trCriticalValue，而 $x$ 是躺在\trFiber{}$(\Psi^k)^{-1}(w)$ 中的那个唯一\trCriticalPoint，那么下列附加条件成立：$(D^2\Psi^k)_x$ 限制到 $(D\varpi^k)_x$ 的核上仍然\trNondegenerate；函数
\begin{equation*}
(\Psi^k)^{-1}(w)\setminus\{x\}\xrightarrow{\ \varpi^k\ }\{w\}\times\CC=\CC
\end{equation*}
具有\trNondegenerate{}的\trCriticalPoint；并且任何两个这种点的取值两两不同，也都不同于 $\varpi^k(x)$ 本身。
\end{itemize}
这些性质\trGeneric{}地成立，这是经典的事实，但为方便起见，我们仍然给出证明的梗概。第一部分是 Bertini--Sard 定理的直截了当的应用。第二部分可以通过下面的标准论证化归为同一定理：考虑\trPullback{}丛 $T^{*}X^k\to\CC^{N-k}\times X^k$ 的、由 $(\lambda,x)\mapsto d\bigl(\sum_{i\ge k}\lambda_i x_i\bigr)$ 给出的那个\trSection。它总与\trZeroSection\trTransverse，于是它的零点集 $Z\subset\CC^{N-k}\times X^k$ 是\trSmooth{}的。如果取\trProjection{}$Z\to\CC^{N-k}$ 的一个\trRegularValue{}$\lambda$，那么函数 $\sum_i\lambda_i x_i$ 只有\trNondegenerate{}的零点。显然，任何一个这样的函数都可以经一次\trLinear{}坐标变换变成 $x_k$。这表明，给定任何 $\bar{X}$，我们可以用 $\mathrm{PGL}_{N+1}$ 中一个贴近恒同的元素来改变它的\trEmbedding，使得变换后的簇上的函数 $\Psi^k=x_k$ 只有\trNondegenerate{}的\trCriticalPoint。对\trEmbedding{}再作一步\trPerturbation，就会使诸\trCriticalPoint{}落在不同的\trFiber{}之中。同样的论证适用于 $\varpi^k$。读者想必已经注意到，这些条件正是\cref{sub:bifibrations}中那些条件的代数几何对应物。其蕴含也与之类似：如果 $x$ 是 $\Psi^k$ 的一个\trCriticalPoint，取值为 $w$，那么 $\mathrm{Critv}(\varpi^k)$ 与直线 $\{w\}\times\CC$ 在 $\varpi^k(x)$ 处相切，而在其余各处都与该直线\trTransverse。定义 $C^k=\mathrm{Critv}(\Psi^k)$（$0\le k\le n$）与 $F^k=\mathrm{Fakev}(\Psi^k,\varpi^k)$（$0\le k\le n-1$），一如\cref{sub:bifibrations}中那样。为简化记号，令 $F^n=\varnothing$。注意，按假设 $0\notin C^k\cup F^k$。

取定 $0\le k\le n$，并写 $C^k=\{c^k_j\}$。为每个 $j$ 选一条\trEmbedding{}路径 $\gamma^k_j\colon[0;1]\to\CC$，满足 $\gamma^k_j(0)=0$、$\gamma^k_j(1)=c^k_j$、对 $t<1$ 有 $\gamma^k_j(t)\notin C^k$，并且对一切 $t$ 有 $\gamma^k_j(t)\notin F^k$。这些路径应按\cref{sub:distinguished-bases}中的意义构成一个\trBasis。在最低维的情形 $k=n$，这些正是从 $0$ 通到\trBranchedCover{}$\Psi^n\colon X^n\to\CC$ 的诸\trBranchPoint{}之一的路径；可以给每条路径配一个无序的点偶 $\mu^{n+1}_j\subset X^{n+1}$，它标出\trCovering{}中沿该路径相会合的那两叶。在更高维数，则按\cref{sub:vanishing-as-matching}行事：追踪子集 $\mathrm{Critv}(\varpi^k)\cap(\{\gamma^k_j(t)\}\times\CC)\subset\CC$ 对 $t\in[0;1)$ 如何移动。这产生一条\trEmbedding{}路径 $\mu^{k+1}_j\colon[-1;1]\to\CC$，满足 $(\mu^{k+1}_j)^{-1}(C^{k+1})=\{-1;1\}$，在相差\trIsotopy{}与\trOrientationReversing{}的意义下唯一。这一过程的产出是下列信息，即 $X$ 的\emph{\trPicard--\trLefschetz{}数据}（Picard--Lefschetz data）：
\begin{itemize}
\item 有限集 $X^{n+1}$，连同子集族 $\{\mu^{n+1}_j\}$；
\item 对每个 $0\le k\le n$：在带\trMarkedPoint{}$\{0\}\cup C^k\cup F^k$ 的平面 $\CC$ 中，路径族 $\{\gamma^k_j\}$，以及（设 $k>0$）$\{\mu^k_j\}$；
\item 在 $H^1(\bar{X})\neq 0$ 的情形，一份附加的\trHomotopy{}数据，马上就会具体给出。
\end{itemize}
注意，这些意在成为\trTopology{}图像（可以移动那些\trMarkedPoint，并为相关的路径配上一个适当的\trIsotopy，而不以任何本质的方式改变数据）。于是，每条路径都可以通过指明如何经\trPunctured{}平面的一个\trDiffeomorphism{}从一个标准路径得到它来编码，这主要给我们留下\trBraidGroup{}中的一些字（这是 Moishezon 的\trBraid{}因子分解技术的一个弱版本）。

让我们看看这些数据含有多少关于 $X$ 的信息。首先，分歧数据使人能够从\trTopology{}上重构 $\Psi^n\colon X^n\to\CC$。为了给 $X^n$ 得到一个\trWellDefined{}的\trFukayaCategory，还需要知道 $\mathcal{K}_{X^n}$ 的\trTrivialization，至少从\trTopology{}上知道（作为在\trHomotopy{}意义下考虑的 $C^{\infty}$-\trTrivialization）。按构造，这一\trTrivialization{}来自一个处处非零的复 1-\trForm，它在无穷远处的每一点有 $(d-n)$ 阶极点。为了指明它的\trHomotopy{}类，我们应当为一张张成 $H_1(X^n)$ 的\trImmersion{}曲线族写下“\trRotationNumber{}”（即\trMaslovIndex）。如果一条曲线是 $\Psi^{n-1}$ 的\trVanishingCycle，我们已经知道该\trRotationNumber{}为零（因为这条曲线容许一个\trGrading）。假设 $H^1(\bar{X})=0$。那么由弱 \trLefschetz{}定理 $H^1(\bar{X}^{n-1})=0$，而诸\trVanishingCycle{}张成 $H^1(\bar{X}^n)$，这意味着 $\mathcal{K}_{X^n}$ 的\trTrivialization{}是唯一确定的。否则，可能还需要记住一些附加的\trRotationNumber，这就是上面提到的最后一份数据。
\begin{remark}\label{thm:branch-point-alternative}
另一种做法（稍微有悖于总体的哲学）是假设对\trBranchPoint{}集 $C^n$ 的精确了解。在那种情形，$X^n$ 在相差\trHolomorphic{}\trIsomorphism{}的意义下确定，并且它带着一个唯一的、无零点的\trHolomorphic{}1-\trForm，它在无穷远各点周围具有所需的极点阶。
\end{remark}
在下一个层次上，$X^n$ 中有唯一的简单闭曲线，它们在 $\Psi^n$ 下的像是 $\mu^n_j$。此外，这些曲线没有一条是零\trHomology{}的（因为它们是作为\trVanishingCycle{}而产生的，从而对\trForm{}的一个适当选取是\trExact{}的）。循着 \trDonaldson{}与 Gompf 的思路，这一曲线族连同相应的路径 $\gamma^{n-1}_j$ 一起，在相差辛\trIsomorphism{}的意义下确定 $\bar{X}^{n-1}$（相对于无穷远\trDivisor，从而在 $X^{n-1}$ 上诱导一个\trExact{}的辛\trIsomorphism）。最后，$\mathcal{K}_{X^n}$ 的\trTrivialization{}以在\trHomotopy{}意义下唯一的方式诱导 $\mathcal{K}_{X^{n-1}}$ 的一个。

对 $k<n$，\trPicard--\trLefschetz{}数据本身是否足以确定 $X^k$ 或 $\bar{X}^k$ 的辛\trIsomorphism{}型，看来是存疑的。这是因为我们记住的是\trMatchingPath，却没有记住构造实际的\trMatchingCycle{}所需要的那些\trIsotopy（经\trCompactification{}把 $X^k$ 变成 $\bar{X}^k$ 时，实际上还有另一个潜在的问题，即无穷远\trDivisor{}被贴附的方式）。人们可以试着通过指明附加的数据来克服这些问题，但这不是我们选择要走的方向。
